%=== Matching ===
\question\MatchQuestion{Infauna}{All benthic organisms that live \emph{in} the sediment.}
\vspace{0.7\baselineskip}
\question\MatchQuestion{Mixed Semidiurnal}{One very low and one moderately low tide on a daily cycle.}
\vspace{0.7\baselineskip}

%=== True False =====

% FIX \question\TFQuestion{False}{Freshwater would extend farther up on the right side of an estuary in the northern hemisphere. Salt water would extend farther down on the left side of the estuary.}

\question\TFQuestion{False}{Strong winds are more likely to decrease GPP in a coastal plain estuary than a bar-built estuary.}

%=== Multiple Choice ===

%BeginMC
\question
An osmoconformer  is \rule{1in}{0.4pt}

\begin{choices}
\choice an estuarine organism whose blood salinity conforms to open ocean salinity. %EndChoice
\choice an organism that maintains a constant blood salinity at any external salinity. %EndChoice
\correctchoice an organism whose blood salinity matches any  external salinity. %EndChoice
\choice an organism that maintains a constant rate of osmoregulation. %EndChoice
\choice any stenohaline organism. %EndChoice
\choice any euryhaline organism. %EndChoice
\end{choices}
%EndMC

%=== Short answer ===

\question[5]
Briefly describe one reason why the size of a tide pool (larger and deeper vs. smaller and shallower) might determine the types of organisms present in the pool.
	\begin{solution}
	Smaller/shallower tide pools are more vulnerable to extreme fluctuations in temperature and salinity, or even drying out. Organisms found in small tide pools will have greater tolerance for the fluctuating conditions.
	\end{solution}
\vspace*{\stretch{1}}


\question[5]
List in order from ocean to land the four zones found in a typical New England salt marsh. You may use common or scientific names but scientific names make a greater impression. 

\AnswerBox{0.1\textheight}{Cordgrass zone, salt hay zone, black rush zone, marsh elder zone.}

\newpage


%% 10 Point or so essays.

\question[10]
Compare and contrast a horizontal hard-bottom community to a vertical hard-bottom community in terms of the organisms likely to be present and the important abiotic and biotic factors that determine the differences.

	\begin{solution}
	Vertical surfaces receive less light than horizontal surfaces, so tends to have less algae. Because less algae is present, it tends to be covered in sessile organisms. (5 points)

	Horizontal surfaces have more algae due to more light. However, sea urchins may graze on the algae, leaving more space for other sessile organisms. Sea urchins not present on vertical surfaces because tube feet can't keep them attached. (5 points)
	\end{solution}
\vspace*{\stretch{2}}
\newpage


\question[10]
Briefly explain how and why spring and neap tides are different from each other. Diagram the relationship between the sun, earth and moon for spring tides and neap tides. You do not have to include the moon phases.
\vspace*{\stretch{3}}

\question[10]
Describe the \emph{specific} biotic and abiotic interactions between Cordgrass (\textit{Spartina alterniflora}) and Salt Hay (\textit{S. patens}) that determine their distribution in a salt marsh. Include in your answer which species is found in the low marsh and which is found in the high marsh. 

\begin{AnswerPage}{0.6\textheight}%
	1. Patens (high marsh) is better competitor for space but it cannot tolerate wet, low O\textsubscript{2} of sediment in low marsh.
	
	2. Alterniflora cannot compete with patens for space but can tolerate wet, low O\textsubscript{2} of sediment in low marsh.
\end{AnswerPage}


\question[10]
Compare and contrast the predator/disturbance and recruitment limitation models proposed to explain high diversity of coral reef fish communities. Explain the roles of competition and post-recruitment processes in your answer.

\AnswerBox{0.4\textheight}{Both models have weak competition. The predation/disturbance model proposes that there is an abundance of new larval recruits but predation or disturbance removes some of the recruits from the population. This keeps the population size of each species well below carrying capacity so competition among individuals is weak. Recruitment limitation suggests that there is a there is a limited number of new larval recruits. Because they are limited in abundance, the populations do not reach carrying capacity so competition is weak. There are no processes that affect the relative abundance of the new larval recruits so post-recruitment modification is absent.}


%% 15 point essays

\question
Answer the next two questions with regard to a bar-built estuary. Write precise answers and use proper terminology throughout. 
	\begin{parts}
	\part[8]
	Describe the distribution of sediment types in the estuary, and explain why the sediments have this distribution.
		\begin{solution}
		(5 pts) Larger sand grains will be closer to the river or in areas where velocity is faster because they are heavier and settle out sooner. Immediately behind the bar will be sand again because wave action can stir up the fine sediments.
		\vspace{\baselineskip}

		(3 pts) Silt and clay (mud) are farther from the river in the quieter areas.
		\end{solution}
	\vspace*{\stretch{1.2}}

	\part[8]
	Describe the distribution of suspension and deposit feeders in the estuary, and explain why each appears where it does in the estuary. 
		\begin{solution}
		(4 pts) Suspension feeders are most common in sandy areas. They can't live in muddy areas because the silt and clay will clog their feeding apparatus. 
		\vspace{\baselineskip}

		(4 pts) Deposit feeders are found with mud. They ingest the mud and digest the organic material. The organic material settles out at the same rate as silt and clay so is associated with those sediment types.
		\end{solution}
	\vspace*{\stretch{1}}
	\end{parts}

\question[15]
Compare and contrast salt marshes to mangrove forests. Discuss similarities and differences.  Include examples of the types and relative ages (larvae, juveniles, adults) of representative organisms likely to be present in different parts of each type of estuary.




